\begingroup
\section{Node A0: strict transverse side activation}
\label{module:A0}
\noindent\textit{Audited source: \nolinkurl{A0_side_activation.tex}.}
\subsection{A sharp one-dimensional tent inequality}

For \(L>0\) and \(x\in(0,L)\), let
\begin{equation}\label{A0:eq:green}
 K_x(s)=
 \begin{cases}
  \displaystyle\frac{s(L-x)}{L},&0\le s\le x,\\[2mm]
  \displaystyle\frac{x(L-s)}{L},&x\le s\le L.
 \end{cases}
\end{equation}
Thus \(K_x\) is the Dirichlet Green kernel of \(-d^2/ds^2\), viewed as a
function of \(s\), and the unit-height triangular tent is
\begin{equation}\label{A0:eq:tent}
 h_x(s)=\frac{L}{x(L-x)}K_x(s).
\end{equation}

\begin{lemma}[Log-concave density dominates its constant rearrangement]
\label{A0:lem:tent}
Let \(q:(0,L)\to(0,\infty)\) be log-concave and integrable, with
\begin{equation}\label{A0:eq:endpoint}
 \lim_{s\downarrow0}q(s)=\lim_{s\uparrow L}q(s)=0.
\end{equation}
Put
\begin{equation}
 \mu=\frac1L\int_0^Lq(s)\,ds.
\end{equation}
If
\begin{equation}\label{A0:eq:barycenter}
 \int_0^L\left(s-\frac L2\right)q(s)\,ds=0,
\end{equation}
then, for every \(x\in(0,L)\),
\begin{equation}\label{A0:eq:tent-positive}
 \boxed{\int_0^Lh_x(s)\{q(s)-\mu\}\,ds>0.}
\end{equation}
\end{lemma}

\begin{proof}
Set \(g=q-\mu\) and
\begin{equation}
 G(y)=\int_0^yg(s)\,ds,\qquad
 H(y)=\int_0^yG(r)\,dr.
\end{equation}
Mass and barycenter matching give
\begin{equation}\label{A0:eq:moments}
 \int_0^Lg=0,\qquad \int_0^Lsg(s)\,ds=0,
\end{equation}
and therefore
\begin{equation}\label{A0:eq:Hend}
 H(L)=\int_0^L(L-s)g(s)\,ds=0.
\end{equation}

Log-concavity says that every superlevel set of \(q\) is an interval.
In view of \eqref{A0:eq:endpoint} and \(\int q=\mu L\), the set
\(\{q>\mu\}\) is a nonempty interval \((a,b)\Subset(0,L)\), up to
irrelevant closed plateau endpoints.  Thus \(g\) has the sign pattern
\[
 -,\quad +,\quad -.
\]
Consequently \(G\) first decreases, then increases, and finally decreases
to \(G(L)=0\).  On the last interval it must be nonnegative, because it is
decreasing to zero.  Hence \(G\) changes sign exactly once, from negative
to positive.  It follows that \(H\) first decreases and then increases.
Together with \(H(0)=H(L)=0\), this gives
\begin{equation}\label{A0:eq:Hnegative}
 H(y)<0,\qquad 0<y<L.
\end{equation}
Strictness follows because \(q\) cannot equal the positive constant
\(\mu\) while satisfying \eqref{A0:eq:endpoint}.

Now set
\[
 F(x)=\int_0^LK_x(s)g(s)\,ds.
\]
The Green-kernel identity gives
\[
 -F''=g,\qquad F(0)=F(L)=0.
\]
Moreover \(F'(0)=L^{-1}\int_0^L(L-s)g(s)\,ds=0\) by
\eqref{A0:eq:moments}.  Hence \(F'=-G\) and \(F=-H\).  Equation
\eqref{A0:eq:Hnegative} yields \(F(x)>0\).  Multiplication by the positive
factor in \eqref{A0:eq:tent} proves \eqref{A0:eq:tent-positive}.
\end{proof}

\begin{remark}
The proof is a convex-order statement, but no external convex-order
theorem is needed.  The only special input is that a log-concave density
crosses its average on one interval; the two moment equalities then force
the sign of every tent pairing.
\end{remark}

\subsection{Log-concavity of the side flux}

Let \(P\) be a bounded convex polygon and let \(u>0\) be its first
Dirichlet eigenfunction:
\begin{equation}
 -\Delta u=\lambda u\quad\hbox{in }P,\qquad
 u=0\quad\hbox{on }\partial P.
\end{equation}
Fix an open side \(S\), identify it with \(0<s<L\), and let
\(\nu_{\rm in}\) be its inward unit normal.  Define
\begin{equation}\label{A0:eq:p}
 p(s)=\partial_{\nu_{\rm in}}u(s)>0.
\end{equation}

\begin{lemma}[The boundary flux is log-concave on a side]
\label{A0:lem:flux-logconcave}
The function \(p\) in \eqref{A0:eq:p} is log-concave on \((0,L)\).
Moreover,
\begin{equation}\label{A0:eq:p-endpoints}
 p(s)\longrightarrow0
 \quad\hbox{as }s\downarrow0\hbox{ or }s\uparrow L.
\end{equation}
Consequently \(p^2\) satisfies the hypotheses of
Lemma~\ref{A0:lem:tent}.
\end{lemma}

\begin{proof}
The Brascamp--Lieb theorem~\cite{BL76} gives log-concavity of \(u\) in every bounded
convex domain.  Fix \(s_0,s_1\in(0,L)\) and \(0<\theta<1\).  For all
sufficiently small \(r>0\), the three points
\[
 s_j+r\nu_{\rm in}\quad(j=0,1),\qquad
 ((1-\theta)s_0+\theta s_1)+r\nu_{\rm in}
\]
belong to \(P\).  Log-concavity of \(u\) gives
\begin{multline}
 \frac{u(((1-\theta)s_0+\theta s_1)+r\nu_{\rm in})}{r}\\
 \ge
 \left(\frac{u(s_0+r\nu_{\rm in})}{r}\right)^{1-\theta}
 \left(\frac{u(s_1+r\nu_{\rm in})}{r}\right)^\theta.
\end{multline}
Letting \(r\downarrow0\) proves log-concavity of \(p\).

At a convex vertex of interior angle \(\omega<\pi\), the standard corner
expansion of a Dirichlet solution starts with
\(r^{\pi/\omega}\sin(\pi\vartheta/\omega)\).  Hence its gradient is
\(O(r^{\pi/\omega-1})\), which tends to zero.  This proves
\eqref{A0:eq:p-endpoints}.
\end{proof}

\subsection{Activating one new side}

Use the scale-invariant functional
\begin{equation}\label{A0:eq:Lambda}
 \Lambda(P)=\frac{|P|}{\pi}\lambda_1(P).
\end{equation}
Suppose \(P\) is stationary under all active vertex variations.  On each
side \(S\) of length \(L\), stages 70--72 give the exact identities
\begin{align}
 \frac1L\int_0^Lp(s)^2\,ds&=\frac{\lambda}{|P|}=:\mu,
 \label{A0:eq:mean}\\
 \int_0^L\left(s-\frac L2\right)p(s)^2\,ds&=0.
 \label{A0:eq:first-moment}
\end{align}

Choose \(x\in(0,L)\), insert at that point a collinear marked vertex, and
move it a distance \(\varepsilon>0\) in the outward normal direction,
keeping the two endpoints of \(S\) fixed.  For small \(\varepsilon\) this
replaces \(S\) by two sides of a convex polygon and adds a triangular cap.
At \(\varepsilon=0\), the normal velocity on \(S\) is exactly the tent
\(h_x\) of \eqref{A0:eq:tent}.

\begin{theorem}[Strict side-activation descent]\label{A0:thm:descent}
For every side and every \(x\in(0,L)\), the outward splitting path
\(P_\varepsilon\) satisfies
\begin{equation}\label{A0:eq:strict}
 \boxed{
 \frac d{d\varepsilon}\Lambda(P_\varepsilon)
 \bigg|_{\varepsilon=0+}<0.}
\end{equation}
Thus a stationary convex \(M\)-gon is not stationary, and is not a local
minimum, in the closed class of convex polygons with at most \(M+1\)
active sides.
\end{theorem}

\begin{proof}
Hadamard's formula and the area derivative give, for an outward normal
velocity \(h\),
\begin{align}
 D\Lambda(P)[h]
 &=\frac{\lambda}{\pi}\int_{\partial P}h\,ds
   -\frac{|P|}{\pi}\int_{\partial P}p^2h\,ds\notag\\
 &=\frac{|P|}{\pi}\int_{\partial P}
   \left(\frac{\lambda}{|P|}-p^2\right)h\,ds.
 \label{A0:eq:Hadamard}
\end{align}
For the splitting path, \(h=h_x\) on \(S\) and is zero on the other
sides.  Lemma~\ref{A0:lem:flux-logconcave} applies to \(q=p^2\), while
\eqref{A0:eq:mean}--\eqref{A0:eq:first-moment} are precisely its mass and
barycenter hypotheses.  Lemma~\ref{A0:lem:tent} therefore gives
\[
 \int_0^Lh_x(p^2-\mu)\,ds>0.
\]
Insert this inequality and \(\mu=\lambda/|P|\) into
\eqref{A0:eq:Hadamard} to obtain \eqref{A0:eq:strict}.

The path has one additional active vertex for every sufficiently small
\(\varepsilon>0\).  Since its scale-invariant value decreases to first
order, the final assertion follows.
\end{proof}

\begin{corollary}[No minimizing side loss]\label{A0:cor:no-loss}
Let \(P\) minimize \(\Lambda\) in the class of convex polygons with at
most \(N\) sides.  If \(P\) has at least three sides and is stationary in
its active stratum, then it has exactly \(N\) active sides.
\end{corollary}

\begin{proof}
If \(P\) had \(M<N\) active sides, Theorem~\ref{A0:thm:descent} would produce
a convex \((M+1)\)-gon with strictly smaller \(\Lambda\), contradicting
minimality.
\end{proof}
\endgroup
